The main experiment uses the earlier version's DQN settings. To check that the encoder effect and the reward comparison are not specific to them, we retrain the monitorability-aware agents with both encoders, and outcome-only and naive agents with the invariant encoder, under three alternatives, five seeds each: learning rate $3{\cdot}10^{-4}$; hard target updates every 1{,}000 steps; and Polyak-averaged targets ($\tau = 0.005$)~\citep{lillicrap2016continuous} (Table~\ref{tab:sensitivity}). None of them removes the churn. In every setting the encoder gap reproduces: on lengths 10--14 the invariant encoder beats the absolute one by $+$\val{sens/ma_inv_lr3e-4-ma_abs_lr3e-4/heldout_xlong/success}, $+$\val{sens/ma_inv_sync1000-ma_abs_sync1000/heldout_xlong/success} and $+$\val{sens/ma_inv_polyak-ma_abs_polyak/heldout_xlong/success} respectively, in \val{sens/ma_inv_lr3e-4-ma_abs_lr3e-4/heldout_xlong/success/npos}/5, \val{sens/ma_inv_sync1000-ma_abs_sync1000/heldout_xlong/success/npos}/5 and \val{sens/ma_inv_polyak-ma_abs_polyak/heldout_xlong/success/npos}/5 seeds. With five seeds, the smallest attainable sign-flip $p$-value is $0.0625$.

The reward comparison also holds (Table~\ref{tab:sensitivity_rewards}), including with twice the training (40{,}000 episodes, with the exploration schedule stretched to match; outcome-only and monitorability-aware agents only). In every setting, late checkpoints of the monitorability-aware agents collapse far less often (\val{sensitivity/ma_inv_lr3e-4/collapse_rate/pct}, \val{sensitivity/ma_inv_sync1000/collapse_rate/pct} and \val{sensitivity/ma_inv_polyak/collapse_rate/pct}) than those of outcome-only agents (\val{sensitivity/outcome_inv_lr3e-4/collapse_rate/pct}, \val{sensitivity/outcome_inv_sync1000/collapse_rate/pct} and \val{sensitivity/outcome_inv_polyak/collapse_rate/pct}). The monitorability-aware agents beat outcome-only at lengths 7--9 by $+$\val{sens/ma_inv_lr3e-4-outcome_inv_lr3e-4/heldout_long/success}, $+$\val{sens/ma_inv_sync1000-outcome_inv_sync1000/heldout_long/success} and $+$\val{sens/ma_inv_polyak-outcome_inv_polyak/heldout_long/success}, in \val{sens/ma_inv_lr3e-4-outcome_inv_lr3e-4/heldout_long/success/npos}/5, \val{sens/ma_inv_sync1000-outcome_inv_sync1000/heldout_long/success/npos}/5 and \val{sens/ma_inv_polyak-outcome_inv_polyak/heldout_long/success/npos}/5 seeds. Naive process reward again collapses less often than outcome-only (\val{sensitivity/naive_inv_lr3e-4/collapse_rate/pct}, \val{sensitivity/naive_inv_sync1000/collapse_rate/pct} and \val{sensitivity/naive_inv_polyak/collapse_rate/pct}), but whether it does better on held-out lengths depends on the learner: at 7--9 it is ahead by \val{sens/naive_inv_lr3e-4-outcome_inv_lr3e-4/heldout_long/success} with the lower learning rate (\val{sens/naive_inv_lr3e-4-outcome_inv_lr3e-4/heldout_long/success/npos}/5 seeds) and behind with Polyak averaging (a difference of \val{sens/naive_inv_polyak-outcome_inv_polyak/heldout_long/success}, \val{sens/naive_inv_polyak-outcome_inv_polyak/heldout_long/success/npos}/5 seeds ahead).

\begin{table}[t]
\centering
\small
\setlength{\tabcolsep}{3pt}
\caption{The reward comparison (invariant observations) under each learner setting: the fraction of late checkpoints (after 5{,}000 episodes) that collapse to training-length success $\le 0.35$, and selected-checkpoint success at lengths 7--9 with its 95\% $t$-interval. Default: 10 seeds; alternatives: 5 seeds.}
\label{tab:sensitivity_rewards}
\resizebox{\linewidth}{!}{\begin{tabular}{lcccccc}
\toprule
Learner & Collapsed: outcome & naive & MA & Success 7--9: outcome & naive & MA \\
\midrule
default (main experiment) & 85\% & 62\% & 21\% & 0.62\,{\scriptsize[0.48, 0.75]} & 0.76\,{\scriptsize[0.64, 0.87]} & 0.86\,{\scriptsize[0.83, 0.89]} \\
lr $3{\cdot}10^{-4}$ & 81\% & 49\% & 19\% & 0.48\,{\scriptsize[0.29, 0.68]} & 0.64\,{\scriptsize[0.43, 0.84]} & 0.82\,{\scriptsize[0.75, 0.89]} \\
hard sync 1000 & 96\% & 75\% & 17\% & 0.40\,{\scriptsize[0.18, 0.62]} & 0.58\,{\scriptsize[0.41, 0.75]} & 0.82\,{\scriptsize[0.74, 0.91]} \\
Polyak $\tau{=}0.005$ & 87\% & 80\% & 19\% & 0.52\,{\scriptsize[0.18, 0.86]} & 0.36\,{\scriptsize[0.26, 0.46]} & 0.86\,{\scriptsize[0.84, 0.88]} \\
40k episodes & 83\% & -- & 22\% & 0.53\,{\scriptsize[0.32, 0.74]} & -- & 0.83\,{\scriptsize[0.75, 0.90]} \\
\bottomrule
\end{tabular}
}
\end{table}

\begin{table}[t]
\centering
\small
\setlength{\tabcolsep}{3pt}
\caption{Sensitivity of the encoder comparison to DQN hyperparameters (monitorability-aware reward; mean over 5 seeds, 95\% $t$-interval). Main-experiment values with the default settings are in Table~\ref{tab:main}.}
\label{tab:sensitivity}
\resizebox{\linewidth}{!}{\begin{tabular}{lccccc}
\toprule
Variant & Success 3--6 & Success 7--9 & Success 10--14 & $P(V\mid c{=}1)$ & $P(V\mid c{=}0)$ \\
\midrule
abs, lr $3{\cdot}10^{-4}$ & 0.72\,{\scriptsize[0.62, 0.81]} & 0.24\,{\scriptsize[0.12, 0.35]} & 0.03\,{\scriptsize[0.00, 0.06]} & 0.76 & 0.11 \\
inv, lr $3{\cdot}10^{-4}$ & 0.88\,{\scriptsize[0.87, 0.90]} & 0.82\,{\scriptsize[0.78, 0.86]} & 0.67\,{\scriptsize[0.56, 0.75]} & 0.95 & 0.11 \\
abs, hard sync 1000 & 0.80\,{\scriptsize[0.78, 0.82]} & 0.29\,{\scriptsize[0.22, 0.38]} & 0.02\,{\scriptsize[0.00, 0.04]} & 0.70 & 0.20 \\
inv, hard sync 1000 & 0.90\,{\scriptsize[0.89, 0.91]} & 0.82\,{\scriptsize[0.76, 0.87]} & 0.63\,{\scriptsize[0.54, 0.70]} & 0.95 & 0.18 \\
abs, Polyak $\tau{=}0.005$ & 0.76\,{\scriptsize[0.74, 0.78]} & 0.26\,{\scriptsize[0.19, 0.33]} & 0.04\,{\scriptsize[0.00, 0.08]} & 0.66 & 0.18 \\
inv, Polyak $\tau{=}0.005$ & 0.90\,{\scriptsize[0.89, 0.91]} & 0.86\,{\scriptsize[0.85, 0.87]} & 0.75\,{\scriptsize[0.70, 0.80]} & 0.83 & 0.13 \\
\bottomrule
\end{tabular}
}
\end{table}
